\section{Thảo luận kết quả so với mục tiêu}

Đối chiếu mục tiêu cụ thể Chương~1:

\begin{itemize}
    \item Corpus \soLuat{} văn bản, cắt Điều, metadata và viện dẫn: đạt ở mức
    nhập và hiển thị; chưa đạt lịch sử hiệu lực.
    \item Pipeline bảy bước, SSE: đạt trên cấu hình sản phẩm (rerank tắt).
    \item Tra cứu cây và không dấu: đạt, có test FTS.
    \item Soát hợp đồng PDF/DOCX/TXT nền: đạt với file có lớp chữ; không OCR.
    \item Lịch \soRule{} rule theo hồ sơ: đạt; hạn là công thức trong mã.
    \item JWT, vai trò, cách ly: đạt, có test 404.
    \item Docker Compose và \soTest{} test: đạt trên môi trường báo cáo.
\end{itemize}

Mục tiêu tổng quát ``mọi căn cứ mở được Điều'' đạt với điều kiện corpus đã nạp
và citation đi qua hậu kiểm. Điều kiện đó không biến câu chữ thành tư vấn pháp
lý có chữ ký.

\section{Bài học kỹ thuật}

Ba chỗ dễ sai khi dựng hệ thống kiểu này, đã xử lý trong mã:

\begin{enumerate}
    \item Thứ tự tin nhắn không tin \texttt{created\_at} trong cùng transaction.
    \item Compose không đọc \texttt{backend/.env}.
    \item Cùng origin lúc production làm CORS ``biến mất'' --- đừng copypaste
    \texttt{*} từ máy dev sang máy thật.
\end{enumerate}

Ba chỗ còn mở: hiệu lực văn bản, OCR, đo F2 nội bộ.

\section{Kết luận chương}

Đồ án cho thấy RAG pháp lý tiếng Việt vận hành được trong ứng dụng web đa người
dùng, với hàng rào dẫn chiếu và lịch rule-based, trên ngăn xếp sinh viên triển
khai được. Phần chưa làm được đã liệt kê để hướng dẫn viên và hội đồng khỏi
hiểu nhầm phạm vi.
